\begin{lstlisting}[
    caption={$k$-NN search using Milvus HNSW approach.},
    label={lst:milvus-hnsw}
]
import numpy as np
from pymilvus import FieldSchema, DataType, CollectionSchema, Collection, list_collections

collection_name = "CIFAR10"

# Drop the collection if it already exists to ensure a clean start
if collection_name in list_collections():
    Collection(collection_name).drop()

# Define and create the collection schema
# 'index' serves as the primary key, 'embedding' stores the image vectors
collection = Collection(
    collection_name,
    CollectionSchema([
        FieldSchema(name="index", dtype=DataType.INT64, is_primary=True),
        FieldSchema(name="embedding", dtype=DataType.FLOAT_VECTOR, dim=d)
    ])
)

# Insert embeddings with their indices and flush to persist data
collection.insert([np.arange(len(embeddings)), embeddings])
collection.flush()

# Create an HNSW index on the embedding field
index_params = {"index_type": "HNSW", "metric_type": "L2", "params": {"M": 32, "efConstruction": 40}}
collection.create_index(field_name="embedding", index_params=index_params)

# Load the collection into memory to enable search
collection.load()

# Define search parameters and perform top-k nearest neighbor search
search_params = {"metric_type": "L2", "params": {"ef": 50}}
results = collection.search([query_vector], "embedding", search_params, limit=k)

top_ids = [res.index for res in results[0]]
print("Top-k ids:", top_ids)
print("Top-k labels:", labels[top_ids].tolist())
\end{lstlisting}